\chapter{Your first appendix}
\section{The title of the first section}
The appendices work exactly the same way as chapters, they are numbered with letters rather than numbers though.


NEED Coefficients A-K evaluated in the appendix and referenced in chapter 3.

% ======================================================================
%  Appendix  --  Fitting of the power-law exponents
%  Referenced from Chapter 6 as \ref{app:exponents}
% ======================================================================

\section{Fitting of the power-law exponents}
\label{app:exponents}

Every sensitivity in Chapter~\ref{chap:parametric} is reported as the
exponent $n$ of a power law $y \propto x^{\,n}$ fitted across the swept
range. An exponent is used in preference to a percentage change because it
does not depend on how wide the range happens to be. This appendix sets out
how each one was obtained and lists them all in one place.

\subsection*{Method}

A power law $y = a x^{\,n}$ becomes a straight line under logarithms,
\begin{equation}
\label{eq:loglin}
\ln y = \ln a + n \ln x ,
\end{equation}
so $n$ is the gradient of $\ln y$ against $\ln x$ and is found by ordinary
least squares. Writing $x_i = \ln(\text{input})$ and
$y_i = \ln(\text{output})$ for the $N$ points of a sweep,
\begin{equation}
\label{eq:lsq}
n = \frac{N \sum x_i y_i - \sum x_i \sum y_i}
         {N \sum x_i^2 - \left(\sum x_i\right)^2},
\qquad
\ln a = \frac{\sum y_i - n \sum x_i}{N} .
\end{equation}
The quality of the fit is reported as the coefficient of determination,
\begin{equation}
\label{eq:r2}
R^2 = 1 - \frac{\sum (y_i - \hat y_i)^2}{\sum (y_i - \bar y)^2},
\end{equation}
evaluated on the fitted values $\hat y_i = a x_i^{\,n}$ in physical units
rather than on the logarithms, so that it measures agreement with the
quantity actually of interest.

\subsection*{Worked example: capacity against bar diameter}

The six cases of Table~\ref{tab:sweepdb} give the following.

\begin{table}[H]
\centering
\caption{Least-squares working for $F_{II} \propto d_b^{\,n}$, from the
         bar-diameter sweep at $L = \SI{150}{\milli\metre}$.}
\label{tab:lsqwork}
\small
\begin{tabular*}{\linewidth}{@{\extracolsep{\fill}}cccccc@{}}
\hline
$d_b$ (mm) & $F_{II}$ (kN) & $x = \ln d_b$ & $y = \ln F_{II}$
           & $xy$ & $x^2$ \\
\hline
8  & 8.82560  & 2.079442 & 2.177657 & 4.528310 & 4.324077 \\
10 & 10.76008 & 2.302585 & 2.375843 & 5.470581 & 5.301898 \\
12 & 12.68776 & 2.484907 & 2.540638 & 6.313248 & 6.174761 \\
16 & 16.44306 & 2.772589 & 2.799904 & 7.762981 & 7.687248 \\
20 & 19.98282 & 2.995732 & 2.994873 & 8.971837 & 8.974412 \\
24 & 23.25610 & 3.178054 & 3.146567 & 9.999961 & 10.100026 \\
\hline
$\sum$ & & 15.813308 & 16.035481 & 43.046917 & 42.562423 \\
\hline
\end{tabular*}
\end{table}

Substituting into \eqref{eq:lsq} with $N = 6$,
\begin{align}
n &= \frac{6(43.046917) - (15.813308)(16.035481)}
          {6(42.562423) - (15.813308)^2}
   = \frac{4.707497}{5.313822}
   = 0.8859 , \\[4pt]
\ln a &= \frac{16.035481 - 0.8859 \times 15.813308}{6} = 0.337754
      \quad\Longrightarrow\quad a = 1.4018 ,
\end{align}
so that
\begin{equation}
\label{eq:dbfit}
F_{II} = 1.402\, d_b^{0.886}\ \si{\kilo\newton},
\qquad R^2 = 0.99975 ,
\end{equation}
with $d_b$ in millimetres. The residuals are given in
Table~\ref{tab:lsqresid}; none exceeds \SI{0.7}{\percent}.

\begin{table}[H]
\centering
\caption{Residuals of the fit \eqref{eq:dbfit}.}
\label{tab:lsqresid}
\small
\begin{tabular*}{\linewidth}{@{\extracolsep{\fill}}cccc@{}}
\hline
$d_b$ (mm) & measured $F_{II}$ (kN) & fitted (kN) & error (\%) \\
\hline
8  & 8.826  & 8.846  & $+0.23$ \\
10 & 10.760 & 10.779 & $+0.18$ \\
12 & 12.688 & 12.669 & $-0.15$ \\
16 & 16.443 & 16.346 & $-0.59$ \\
20 & 19.983 & 19.919 & $-0.32$ \\
24 & 23.256 & 23.410 & $+0.66$ \\
\hline
\end{tabular*}
\end{table}

\subsection*{Second worked example: capacity against cover, with exclusions}

The cover sweep differs from the one above in that two of its eight cases
are excluded before fitting. The reason is set out in
Section~\ref{sec:geomcover}; the arithmetic is as follows.

The Area Method averages over a half-disc of radius
$R_{avg} = 1.32\,l_{char} = \SI{26.4}{\milli\metre}$, which at
$l_{item} = \SI{2}{\milli\metre}$ requires
$\lceil 26.4/2 \rceil = 14$ radial items. The number that actually fit
between the bar surface and the outer face is
$\lfloor (r_e - d_b/2)/l_{item} \rfloor$, giving
Table~\ref{tab:clipcheck}.

\begin{table}[H]
\centering
\caption{Whether the averaging disc fits inside the concrete. Fourteen
         radial items are required.}
\label{tab:clipcheck}
\small
\begin{tabular*}{\linewidth}{@{\extracolsep{\fill}}ccccc@{}}
\hline
$r_e$ (mm) & $c = r_e - d_b/2$ (mm) & $c/d_b$ & items that fit & status \\
\hline
25  & 19  & 1.583 & 9  & clipped \\
30  & 24  & 2.000 & 12 & clipped \\
40  & 34  & 2.833 & 17 & admissible \\
50  & 44  & 3.667 & 22 & admissible \\
60  & 54  & 4.500 & 27 & admissible \\
75  & 69  & 5.750 & 34 & admissible \\
90  & 84  & 7.000 & 42 & admissible \\
110 & 104 & 8.667 & 52 & admissible \\
\hline
\end{tabular*}
\end{table}

Applying \eqref{eq:lsq} to the six admissible points, with
$x = \ln(c/d_b)$ and $y = \ln F_{II}$:

\begin{table}[H]
\centering
\caption{Least-squares working for $F_{II} \propto (c/d_b)^{\,n}$, six
         admissible points.}
\label{tab:lsqcov}
\small
\begin{tabular*}{\linewidth}{@{\extracolsep{\fill}}cccccc@{}}
\hline
$c/d_b$ & $F_{II}$ (kN) & $x$ & $y$ & $xy$ & $x^2$ \\
\hline
2.833 & 11.07883 & 1.041454 & 2.405036 & 2.504734 & 1.084626 \\
3.667 & 11.28173 & 1.299283 & 2.423185 & 3.148403 & 1.688136 \\
4.500 & 11.76825 & 1.504077 & 2.465405 & 3.708160 & 2.262249 \\
5.750 & 12.68776 & 1.749200 & 2.540638 & 4.444083 & 3.059700 \\
7.000 & 13.69603 & 1.945910 & 2.617106 & 5.092653 & 3.786566 \\
8.667 & 15.09083 & 2.159484 & 2.714087 & 5.861029 & 4.663372 \\
\hline
$\sum$ & & 9.699409 & 15.165457 & 24.759062 & 16.544650 \\
\hline
\end{tabular*}
\end{table}

\begin{align}
n &= \frac{6(24.759062) - (9.699409)(15.165457)}
          {6(16.544650) - (9.699409)^2}
   = \frac{1.458410}{5.189374} = 0.2810 , \\[4pt]
\ln a &= \frac{15.165457 - 0.2810 \times 9.699409}{6} = 2.073260
      \quad\Longrightarrow\quad a = 7.951 ,
\end{align}
so that
\begin{equation}
\label{eq:covfitapp}
F_{II} = 7.951 \left(\frac{c}{d_b}\right)^{0.281}\ \si{\kilo\newton},
\qquad R^2 = 0.9428 .
\end{equation}

Repeating the fit on all eight points, including the two clipped cases,
returns $F_{II} = 9.199\,(c/d_b)^{0.199}$ with $R^2 = 0.8566$. The two
inadmissible points therefore reduce the exponent by \SI{29}{\percent} and
the coefficient of determination from \num{0.94} to \num{0.86}, which is
why they are excluded.

\subsection*{Third worked example: decomposing the concrete-strength
             exponent}

The exponent on concrete strength is not obtained from a single fit. It is
built from three quantities, because $f'_c$ reaches the capacity by two
separate routes: through the strengths $f_t = f_v = 0.5\sqrt{f'_c}$ that
normalise the interaction index, and through the modulus $E_c$ that
determines how stress is distributed. Each route is measured separately
and the two are then combined.

\paragraph{The three fits.}
Applying \eqref{eq:lsq} to the three data sets gives
Table~\ref{tab:fcfits}.

\begin{table}[H]
\centering
\caption{The three fits behind the concrete-strength exponent.}
\label{tab:fcfits}
\small
\begin{tabular*}{\linewidth}{@{\extracolsep{\fill}}llcrc@{}}
\hline
Fit & Source & $N$ & $n$ & $R^2$ \\
\hline
$\tau_u$ against $f'_c$ & Table~\ref{tab:sweepfc}, $E_c$ tied & 7
  & $+0.3901$ & 0.99832 \\
$\tau_u$ against $E_c$  & Table~\ref{tab:sweepEc}, $f'_c$ fixed & 6
  & $-0.3084$ & 0.99991 \\
$E_c$ against $f'_c$    & \textcite{as3600ConcreteStructures2018} Tab.~3.1.2 & 7
  & $+0.3562$ & 0.98034 \\
\hline
\end{tabular*}
\end{table}

\paragraph{Combining them.}
Since $\tau_u$ depends on $f'_c$ through both $f_t$ and $E_c$, the chain
rule in logarithmic variables gives
\begin{equation}
\label{eq:chain}
\frac{\mathrm{d}\ln \tau_u}{\mathrm{d}\ln f'_c}
  = \frac{\partial \ln \tau_u}{\partial \ln f_t}
    \frac{\mathrm{d}\ln f_t}{\mathrm{d}\ln f'_c}
  + \frac{\partial \ln \tau_u}{\partial \ln E_c}
    \frac{\mathrm{d}\ln E_c}{\mathrm{d}\ln f'_c} .
\end{equation}

The first term is known exactly rather than fitted. From
$f_t = 0.5\sqrt{f'_c}$ the second factor is $\tfrac12$, and the first is
unity: the interaction index \eqref{eq:indexrecall} is homogeneous of
degree $-2$ in the strengths, so scaling $f_t$ and $f_v$ together by any
factor scales the failure load by the same factor, giving
$\partial \ln \tau_u / \partial \ln f_t = 1$. The strength route therefore
contributes $1 \times 0.5 = 0.500$.

The second term is the product of two measured quantities from
Table~\ref{tab:fcfits}, $(-0.3084)(0.3562) = -0.1099$. Hence
\begin{equation}
\label{eq:decompapp}
\frac{\mathrm{d}\ln \tau_u}{\mathrm{d}\ln f'_c}
  = 0.500 - 0.110 = 0.390 ,
\end{equation}
against the directly fitted \num{0.3901}. The two agree to four decimal
places.

Read the other way, \eqref{eq:chain} can be solved for the one factor that
was asserted rather than measured:
\begin{equation}
\frac{\partial \ln \tau_u}{\partial \ln f_t}
  = \frac{0.3901 - (-0.3084)(0.3562)}{0.5} = 1.0000 ,
\end{equation}
which confirms the homogeneity argument numerically.

\paragraph{Local exponents.}
The exponent is not quite constant across the range, because the modulus
table flattens at high strength. Table~\ref{tab:fcdrift} applies
\eqref{eq:decompapp} to each interval in turn.

\begin{table}[H]
\centering
\caption{Local exponent on concrete strength, measured against the value
         predicted by \eqref{eq:decompapp} applied interval by interval.}
\label{tab:fcdrift}
\small
\begin{tabular*}{\linewidth}{@{\extracolsep{\fill}}ccccc@{}}
\hline
Interval (MPa) & $n(E_c)$ & predicted $n(\tau_u)$ & measured $n(\tau_u)$
               & difference \\
\hline
20--25 & 0.478 & 0.353 & 0.348 & $-0.004$ \\
25--32 & 0.486 & 0.350 & 0.349 & $-0.002$ \\
32--40 & 0.385 & 0.381 & 0.381 & $-0.001$ \\
40--50 & 0.265 & 0.418 & 0.420 & $+0.002$ \\
50--65 & 0.275 & 0.415 & 0.417 & $+0.002$ \\
65--80 & 0.275 & 0.415 & 0.417 & $+0.002$ \\
\hline
\end{tabular*}
\end{table}

The agreement holds to within \num{0.005} at every interval, so the
decomposition does more than reproduce the average exponent: it predicts
how that exponent varies across the range, from the modulus table alone.

These values use the same two-point formula as
\eqref{eq:localslope}. For the interval $f'_c = \SI{40}{}$ to
\SI{50}{\mega\pascal}, over which $E_c$ rises from \SI{32800}{} to
\SI{34800}{\mega\pascal} and $\tau_u$ from \SI{2.443}{} to
\SI{2.683}{\mega\pascal},
\begin{align}
n(E_c) &= \frac{\ln(34800/32800)}{\ln(50/40)}
        = \frac{0.059189}{0.223144} = 0.265 , \\[2pt]
\text{predicted } n(\tau_u) &= 0.500 + (-0.3084)(0.265) = 0.418 , \\[2pt]
\text{measured } n(\tau_u) &= \frac{\ln(2.683/2.443)}{\ln(50/40)}
        = 0.420 .
\end{align}

\subsection*{Fourth worked example: the shear-to-tensile strength ratio}

The exponent quoted in Section~\ref{sec:matratio} is fitted in the same
way as the others, but it describes a relationship that is not in fact a
power law. Both the fit and the exact expression are given here.

\paragraph{The fitted exponent.}
Applying \eqref{eq:lsq} to the six cases of Table~\ref{tab:sweepratio},
with $x = \ln(f_v/f_t)$ and $y = \ln F_{II}$:

\begin{table}[H]
\centering
\caption{Least-squares working for $F_{II} \propto (f_v/f_t)^{\,n}$.}
\label{tab:lsqfv}
\small
\begin{tabular*}{\linewidth}{@{\extracolsep{\fill}}cccccc@{}}
\hline
$f_v/f_t$ & $F_{II}$ (kN) & $x$ & $y$ & $xy$ & $x^2$ \\
\hline
0.50 & 7.480  & $-0.693147$ & 2.012233 & $-1.394773$ & 0.480453 \\
0.75 & 10.406 & $-0.287682$ & 2.342383 & $-0.673861$ & 0.082761 \\
1.00 & 12.688 & $ 0.000000$ & 2.540657 & $ 0.000000$ & 0.000000 \\
1.25 & 14.415 & $ 0.223144$ & 2.668269 & $ 0.595407$ & 0.049793 \\
1.50 & 15.709 & $ 0.405465$ & 2.754234 & $ 1.116746$ & 0.164402 \\
2.00 & 17.418 & $ 0.693147$ & 2.857504 & $ 1.980671$ & 0.480453 \\
\hline
$\sum$ & & 0.340927 & 15.175279 & 1.624189 & 1.257862 \\
\hline
\end{tabular*}
\end{table}

\begin{equation}
n = \frac{6(1.624189) - (0.340927)(15.175279)}
         {6(1.257862) - (0.340927)^2}
  = \frac{4.571477}{7.430941} = 0.6152 ,
\end{equation}
with $R^2 = 0.967$. That is the weakest fit of any reported in
Chapter~\ref{chap:parametric}, and individual points are in error by up to
\SI{6.5}{\percent} --- an order of magnitude worse than the other sweeps.
The reason is that the underlying relationship is not a power law.

\paragraph{The exact expression.}
Write the averaged stresses at the governing station, evaluated at the
reference load $F_{ref}$, as $\bar\sigma_\theta = a f_t$ and
$\bar\tau_{rz} = b f_t$. Putting $f_v = k f_t$ into
\eqref{eq:indexrecall} gives an index of $a^2 + b^2/k^2$ at that load, and
since the index is homogeneous of degree two in the load,
\begin{equation}
\label{eq:fvexact}
F_{II}(k) = \frac{F_{ref}}{\sqrt{a^2 + b^2/k^2}} .
\end{equation}
The two constants follow from the baseline case. At $k = 1$ the normalised
stresses at failure are $\hat\sigma = 0.612$ and $\hat\tau = 0.791$ with
$F_{II} = \SI{12.688}{\kilo\newton}$, so scaling back to
$F_{ref} = \SI{50}{\kilo\newton}$,
\begin{equation}
a = 0.612 \times \frac{50}{12.688} = 2.412 ,
\qquad
b = 0.791 \times \frac{50}{12.688} = 3.117 .
\end{equation}

\begin{table}[H]
\centering
\caption{The closed form \eqref{eq:fvexact} against the computed values
         and against the power-law fit.}
\label{tab:fvexact}
\small
\begin{tabular*}{\linewidth}{@{\extracolsep{\fill}}crrrrr@{}}
\hline
$f_v/f_t$ & computed & \eqref{eq:fvexact} & error & power law & error \\
          & (kN) & (kN) & (\%) & (kN) & (\%) \\
\hline
0.50 & 7.480  & 7.480  & $ 0.00$ & 7.908  & $+5.7$ \\
0.75 & 10.406 & 10.405 & $-0.01$ & 10.148 & $-2.5$ \\
1.00 & 12.688 & 12.687 & $-0.01$ & 12.113 & $-4.5$ \\
1.25 & 14.415 & 14.413 & $-0.02$ & 13.895 & $-3.6$ \\
1.50 & 15.709 & 15.706 & $-0.02$ & 15.544 & $-1.1$ \\
2.00 & 17.418 & 17.412 & $-0.03$ & 18.554 & $+6.5$ \\
\hline
\end{tabular*}
\end{table}

Equation~\eqref{eq:fvexact} reproduces every case to better than
\SI{0.1}{\percent}. Its two limits are the reason the power law fails. For
$k \to 0$ the shear term dominates and $F_{II} \to F_{ref} k / b$, linear
in $k$; for $k \to \infty$ the shear term vanishes and
$F_{II} \to F_{ref}/a = \SI{20.7}{\kilo\newton}$, a ceiling set by hoop
tension alone. The local slope is therefore
\begin{equation}
\frac{\mathrm{d}\ln F_{II}}{\mathrm{d}\ln k}
  = \frac{b^2/k^2}{a^2 + b^2/k^2},
\end{equation}
which falls from \num{0.870} at $k = 0.5$ to \num{0.295} at $k = 2$. The
fitted \num{0.615} is an average over that range and should not be
extrapolated beyond it.

\subsection*{All fitted exponents}

Table~\ref{tab:allexp} collects every exponent quoted in
Chapter~\ref{chap:parametric}, each obtained by \eqref{eq:lsq} applied to
the corresponding sweep.

\begin{table}[H]
\centering
\caption{Power-law exponents for every factor, fitted by
         \eqref{eq:lsq}. A positive exponent means the output rises with
         the input.}
\label{tab:allexp}
\small
\begin{tabular*}{\linewidth}{@{\extracolsep{\fill}}lcrrr@{}}
\hline
Factor $x$ & $N$ & $n(F_{II})$ & $n(\tau_u)$ & $n(k_s)$ \\
\hline
Embedment length $L$                & 8 & $+0.131$ & $-0.869$ & $+0.064$ \\
Bar diameter $d_b$ (fixed $L$)      & 6 & $+0.886$ & $-0.114$ & $+1.318$ \\
Bar diameter $d_b$ (fixed $L/d_b$)  & 6 & $+0.898$ & $-1.102$ & $+1.318$ \\
Cover ratio $c/d_b$\,$^{\dagger}$   & 6 & $+0.281$ & $+0.281$ & $-0.289$ \\
Rib thickness $t_{rib}$             & 6 & $-0.001$ & $-0.001$ & $-0.001$ \\
Concrete strength $f'_c$            & 7 & $+0.390$ & $+0.390$ & $+0.137$ \\
Concrete modulus $E_c$              & 6 & $-0.308$ & $-0.308$ & $+0.384$ \\
Strength ratio $f_v/f_t$\,$^{\S}$   & 6 & $+0.615$ & $+0.615$ & $\pm0$ \\
Poisson's ratio $\nu_c$             & 5 & $+0.022$ & $+0.022$ & $-0.043$ \\
Rib modulus $E_r/E_c$               & 5 & $-0.008$ & $-0.008$ & $+0.013$ \\
Steel modulus $E_s$                 & 5 & $+0.309$ & $+0.309$ & $+0.615$ \\
\hline
\multicolumn{5}{@{}p{0.95\linewidth}@{}}{\raggedright
  $^{\dagger}$~Fitted to the six cases in which the averaging disc is not
  clipped by the free surface; see Section~\ref{sec:geomcover}. Including
  the two clipped cases would give $+0.199$ with $R^2 = 0.857$.\newline

  $^{\S}$~A local average only; the underlying relationship is not a power
  law. See \eqref{eq:fvexact}.} \\
\hline
\end{tabular*}
\end{table}

Three entries in Table~\ref{tab:allexp} are not independent fits but follow
from the others by identity.

Where the embedment is held fixed, $\tau_u = F_{II}/(\pi d_b L)$ differs
from $F_{II}$ only by a constant, so the two exponents are equal. This is
why the $n(F_{II})$ and $n(\tau_u)$ columns coincide for every factor
except the two that change a length.

For the embedment-length sweep the length itself is the variable, so
\begin{equation}
n(\tau_u) = n(F_{II}) - 1 = 0.131 - 1 = -0.869 ,
\end{equation}
and for the bar-diameter sweep at constant slenderness, where
$L = 10 d_b$, the bonded area grows as $d_b^2$ and
\begin{equation}
n(\tau_u) = n(F_{II}) - 2 = 0.898 - 2 = -1.102 .
\end{equation}

Finally the slip exponent is never fitted independently. Since
$\delta_u = F_u/k_s$,
\begin{equation}
\label{eq:slipexp}
n(\delta_u) = n(F_{II}) - n(k_s) ,
\end{equation}
which for the bar-diameter sweep gives
$0.886 - 1.318 = -0.432$, matching the value obtained by fitting the
computed slips directly to four decimal places.